\section*{Chapter \ref{ch:s2:swap-spread-and-asset-swap-trades} --- Swap-Spread and Asset-Swap Trades}

\begin{solution}{exo:s2:swap-spread-and-asset-swap-trades:1}
$0.05 \times 0.10 = 0.005$ of the notional a year: 50 basis points.
\end{solution}

\begin{solution}{exo:s2:swap-spread-and-asset-swap-trades:2}
$8.5 \times 12 = 102$ basis points of notional.
\end{solution}

\begin{solution}{exo:s2:swap-spread-and-asset-swap-trades:3}
$40 + 5 - 50 = -5$ basis points a year: the charge takes more than the carry.
\end{solution}

\begin{solution}{exo:s2:swap-spread-and-asset-swap-trades:4}
A two-year position ties up balance sheet for little duration, and the bank credit premium in the floating leg still matters at short maturities;
at thirty years the carry per unit of balance sheet is small relative to the duration hedged, and long-dated receivers (pension funds, insurers)
push the swap rate down.
\end{solution}

\begin{solution}{exo:s2:swap-spread-and-asset-swap-trades:5}
Nothing about credit or duration changes on the last days of a quarter; what changes is what banks report. A dip on those days points to the cost of
holding positions on the reporting date.
\end{solution}

\begin{solution}{exo:s2:swap-spread-and-asset-swap-trades:6}
Holders whose own \omterm{def:s2:swap-spread-and-asset-swap-trades:bscost}{balance-sheet cost} is below the carry: funds with cheap term repo, investors who own Treasuries outright, or anyone whose capital is
not measured on the same basis. For them the carry of 38.8 basis points a year is mostly profit.
\end{solution}

\begin{solution}{exo:s2:swap-spread-and-asset-swap-trades:7}
Over the whole ten years it loses 36.8 basis points a year: carry earns 21.7 but marks lose 63.4, because the spread fell from its premium to about
$-39$ as the \omterm{def:s2:swap-spread-and-asset-swap-trades:bscost}{balance-sheet cost} phased in. A carry trade entered before its driver changes pays for the whole repricing in marks.
\end{solution}

\begin{solution}{exo:s2:swap-spread-and-asset-swap-trades:8}
It is not risk-free: the position must be financed in repo and costs balance sheet (50 basis points a year in the chapter's charge), and the spread
can move further negative, marked at a DV01 of 8.5. What looks like an arbitrage is the price of the balance sheet it needs.
\end{solution}

\begin{solution}{pb:s2:swap-spread-and-asset-swap-trades:1}
\begin{enumerate}
\item Swap rate minus government yield; the spread over the floating rate that a swapped bond pays.
\item A premium for the bank credit in the floating leg and for Treasuries' liquidity.
\item Ten-year 58.6 before and 8.3 after; thirty-year 48.3 before and $-21.4$ after.
\item Frictions for holding bonds make negative spreads unsurprising.
\item The return required on capital held against a trade's gross assets, as a running rate.
\item Persistent CIP deviations, strongest at quarter-ends, pointing to regulation.
\item It falls from about 11.5 to about $-38.8$ basis points as the cost phases in.
\item At the last five days of each quarter, twice as deep at year-end: reporting dates.
\item Long the bond in repo, paying fixed; carry, marks, floating over repo, the charge.
\item $-24.3$ and $-74.3$ before; 32.4 and $-17.6$ after.
\item The spread's noise, at a DV01 of 8.5, is large against the carry.
\item 103.3 basis points at year-ends and 43.9 at other quarter-ends.
\item \textbf{Named result.} After the \omterm{def:s2:swap-spread-and-asset-swap-trades:bscost}{balance-sheet cost} is in, the long-swap-spread trade earns 32.4 basis points a year uncharged and loses 17.6
charged at 50 basis points a year.
\item Where the marginal holder, a bank, is indifferent after its charge.
\item Cheaper balance sheet (a change in rules or in who holds Treasuries).
\item By the spread's volatility times DV01 and the balance sheet available.
\item Charge balance sheet, include the regime change, use repo actually available.
\item Asset-swap carry or the long swap spread.
\item Both are repo-financed Treasury positions whose return depends on funding and collateral.
\item The cost of holding Treasuries on a bank's balance sheet.
\end{enumerate}
\end{solution}

\begin{solution}{iq:s2:swap-spread-and-asset-swap-trades:1}
Swap rate minus government yield at the same maturity; negative means the government bond yields more than the swap, which a holder could lock in
only by financing the bond, which costs repo and balance sheet.
\end{solution}

\begin{solution}{iq:s2:swap-spread-and-asset-swap-trades:2}
Holding the bond uses balance sheet that banks charge for, while receiving fixed on the swap uses little; long-dated receivers keep swap rates low.
Jermann's model with bond-holding frictions gives negative spreads.
\end{solution}

\begin{solution}{iq:s2:swap-spread-and-asset-swap-trades:3}
Buy the bond, pay fixed on a swap matching its coupons and receive floating plus a spread: the asset-swap spread. Pay the repo rate to finance the
bond and the \omterm{def:s2:swap-spread-and-asset-swap-trades:bscost}{balance-sheet cost} of holding it.
\end{solution}

\begin{solution}{iq:s2:swap-spread-and-asset-swap-trades:4}
A further widening negative (a rise in \omterm{def:s2:swap-spread-and-asset-swap-trades:bscost}{balance-sheet costs}), a repo spike, a quarter-end dip deeper than usual, a change in the floating index, and
clearing margin increases.
\end{solution}

\begin{solution}{iq:s2:swap-spread-and-asset-swap-trades:5}
Allocate the capital each position requires under the binding ratio, charge it at the firm's hurdle rate daily, and report desks' P\&L after the
charge, with higher rates for positions held over reporting dates.
\end{solution}

\begin{solution}{iq:s2:swap-spread-and-asset-swap-trades:6}
The bond earns its yield and costs repo; the swap receives floating and pays the fixed rate; summing gives yield minus swap rate plus floating minus
repo. If the floating index moves away from repo (as when an interbank rate is replaced by a secured one), the last term changes and with it the
trade's carry.
\end{solution}
